% Research brief for May 2026, Week 4.
% Focus: architecture, hypothesis, evaluation, and paper path.

\begin{frame}
  \titlepage
  \vspace{-1.0em}
  \begin{center}
    \small Research brief: May 2026, Week 4
  \end{center}
\end{frame}

\begin{frame}{Background}
  \begin{itemize}
    \item Long-context LLMs still pay high prefill/KV cost and often fail to use all context effectively.
    \item Prompt summaries and RAG are external memory mechanisms; they do not give the model an internal updateable state.
    \item Full self-modifying weight updates are attractive but too risky as a first milestone.
    \item Research v0 asks a narrower question: can a wrapper learn compact next-state memory while the base model stays frozen?
  \end{itemize}
\end{frame}

\begin{frame}{Research Significance}
  \begin{block}{Core hypothesis}
    A frozen base model plus a learned memory wrapper can approximate full-context behavior at lower context/KV cost.
  \end{block}

  \begin{block}{Architecture claim}
    \[
      m_t = G_\phi(m_{t-1}, c_t), \qquad y_t = F_\theta(q_t, m_t)
    \]
  \end{block}

  \begin{itemize}
    \item This frames memory compression as next-state prediction.
    \item It creates a safer bridge toward future hidden-state or LoRA/weight memory.
  \end{itemize}
\end{frame}

\begin{frame}{Research Scope}
  \begin{itemize}
    \item Base model: frozen 7B/8B open-source model.
    \item Trainable object: wrapper / explicit memory tokens.
    \item First benchmarks: LoCoMo, LongBench, RULER, Needle-in-a-Haystack.
    \item Later ablations: hidden-state memory, LoRA/weight memory, Doc-to-LoRA comparison.
    \item Do not claim self-modifying LLM yet; v0 is learned memory compression.
  \end{itemize}
\end{frame}

\begin{frame}{Evaluation Criteria}
  \begin{table}
    \small
    \begin{tabular}{p{0.30\textwidth} p{0.58\textwidth}}
      \toprule
      Question & Metric \\
      \midrule
      Does memory preserve quality? & Within 5--10 points of full context. \\
      Does it beat cheap compression? & Better than summary at matched token budget. \\
      Does it really compress? & At least 4x token/KV reduction. \\
      Does sequential update forget? & Early evidence retention drop no more than 5 points. \\
      \bottomrule
    \end{tabular}
  \end{table}
\end{frame}

\begin{frame}{Research Deliverable}
  \begin{itemize}
    \item Paper framing: learned memory wrapper as next-state prediction.
    \item Core comparisons: full context, summary, retrieval, fixed/gist tokens, learned memory.
    \item Main result target: memory wrapper beats summary/retrieval at matched budget while approaching full-context quality.
    \item Risk: if memory tokens cannot beat summary, stop before LoRA/weight-memory escalation.
  \end{itemize}
\end{frame}
